\documentclass[12pt]{article}

% Packages
\usepackage[
  paperwidth=8.5in,
  paperheight=11in,
  margin=0.5in
]{geometry}
\usepackage{amsmath, amssymb}
\usepackage{booktabs}
\usepackage{tabularx}

\newcommand{\tacred}{FS-TACRED}
\newcommand{\qwensmall}{Qwen3-4B}
\newcommand{\gemmasmall}{Gemma3-4B}

\begin{document}

\begin{itemize}
    \item The task we are solving requires random luck because we are directly deriving the answer. 
    In other works on APO, the optimized prompts contain a lot of instructions, and without reasoning directly conclude to yes or no is difficult.
    For direct yes or no answers: it requires luck i.e. random search of the prompt. - my hypothesis
    \item "where to focus" / "where is the important thing" - based important selection of prompt parts from parents
    \item take the difference between child and parent as important items that improved the scores
    \item use LLM itself to identify the differences in the prompt between parent and child that caused the improvement 
    \item Evaluating on the full evaluation set can be expensive (i.e. if I use reasoning). "Automatic Prompt Optimization with “Gradient Descent” and Beam Search" - solves this as multi-arm / multi-bandit problem. check that.
    \item 1-shot to K-shot VS showing how to solve as examples
  \end{itemize}

\paragraph{Top approaches}
\begin{itemize}
  \item GEPA merge: Combine different best sub-modules from different systems to be single best systems
  \item Error taxonomy - updates the prompt only once - i think it is inspired by GRPO i.e. doing batch udpate rather than single update
  \item Evoprompt - a + x (b - c)
  \item SPO - no full validation rather they do pairwise self-judge by LLM based on outputs of different prompts
  \item PromptAgent - uses MCTS to explore new prompt (passes rewards to parent nodes)
  \item a mixed of IO and EO approach - from the Teach better or show smarter
  \item My idea: show trace of changes of the prompt during mutation: that will show the differences and the impact
\end{itemize}

\paragraph{What is necessary}
\begin{itemize}
  \item Grouping of errors - best learning 
  \item Merging best subsystems (multi-agent)
  \item Show trace of changes (idea similar to: CMA-ES (Covariance Matrix Adaptation ES) , DE (used in Evoprompt), or Lamarckian \& Baldwinian Evolution (classic framing) Often implemented via local search + GA, also called memetic algorithms)
  \item MCTS - for search 
  \item Examples Optimization (Random and Mutation from pool)
\end{itemize}

* Use MCTS + UCT for search (exploration + exploitation) \\
Step per iteration:
\begin{enumerate}
\item Select a system (one prompt per system)
\item 
\end{enumerate}

% \bibliographystyle{plain}  
% \bibliography{custom}    

\end{document}
